\chapter{Introdução}
\label{cap:introducao}

A crescente transformação digital observada nas últimas décadas tem impulsionado organizações de diversos setores a buscarem processos cada vez mais eficientes para desenvolvimento, entrega e manutenção de software. Nesse contexto, a adoção de práticas DevOps tornou-se um fator estratégico para reduzir o tempo entre a concepção de uma funcionalidade e sua disponibilização ao usuário final, promovendo maior integração entre equipes de desenvolvimento e operações \cite{dora2024,gitlab2024}.

A evolução das tecnologias de conteinerização, impulsionada por plataformas como Docker e Kubernetes, permitiu que organizações migrassem de modelos baseados em servidores estáticos para arquiteturas altamente escaláveis e distribuídas. O Kubernetes consolidou-se como padrão de mercado para orquestração de contêineres, oferecendo mecanismos de escalabilidade automática, balanceamento de carga, recuperação automática de falhas e abstração da infraestrutura subjacente \cite{jenkinskubernetes,spacelift2024}.

Paralelamente, o crescimento da dependência de bibliotecas de terceiros, componentes de código aberto e serviços externos ampliou significativamente a superfície de ataque das aplicações modernas. Ataques direcionados à cadeia de suprimentos de software (\textit{Software Supply Chain}) tornaram-se mais frequentes e sofisticados, demonstrando que mecanismos tradicionais de segurança baseados exclusivamente em perímetros de rede já não são suficientes para proteger ambientes distribuídos e altamente dinâmicos \cite{gama2024}.

Diante desse cenário, abordagens fundamentadas em \textit{Zero Trust Architecture} (ZTA) têm recebido crescente atenção da comunidade acadêmica e da indústria. O modelo Zero Trust parte do princípio de que nenhum componente deve ser considerado confiável por padrão, exigindo verificação contínua de identidade, contexto e conformidade antes da concessão de qualquer privilégio. Complementarmente, tecnologias como SPIFFE e SPIRE possibilitam a emissão e gerenciamento de identidades criptográficas para \textit{workloads} executados em ambientes distribuídos, fortalecendo mecanismos de autenticação e autorização entre serviços \cite{gama2024}.

Além dos desafios relacionados à segurança, organizações enfrentam a necessidade de mensurar objetivamente a eficiência de seus processos de entrega de software. Nesse contexto, as métricas propostas pelo grupo \textit{DevOps Research and Assessment} (DORA) tornaram-se referência mundial para avaliação de desempenho operacional. Métricas como \textit{Deployment Frequency}, \textit{Lead Time for Changes}, \textit{Change Failure Rate} e \textit{Mean Time to Recovery} permitem avaliar simultaneamente velocidade de entrega e estabilidade operacional \cite{dora2024,dx2024}.

Estudos anteriores têm investigado diferentes dimensões do processo moderno de
entrega de software. As pesquisas conduzidas pelo \textit{DevOps Research and Assessment}
(DORA) consolidaram métricas destinadas à avaliação do desempenho da entrega de software,
considerando aspectos relacionados à vazão e à estabilidade das mudanças
\cite{dora2024}. Em outra perspectiva, \citeonline{gama2024} investigou mecanismos de
conformidade contínua de vulnerabilidades e identificou uma limitação nas abordagens
concentradas exclusivamente nas verificações realizadas durante pipelines CI/CD. Segundo
o autor, tais mecanismos não contemplam adequadamente mudanças dinâmicas ocorridas após
a implantação, como o surgimento de novas vulnerabilidades em aplicações já implantadas.

Para tratar essa limitação, \citeonline{gama2024} propõe o emprego de princípios de
\textit{Zero Trust} e identidades de \textit{workloads} providas pelo SPIRE, permitindo
a reavaliação contínua da postura de segurança e o isolamento de componentes que deixem
de atender aos requisitos mínimos de conformidade. Esse problema também é retomado por
\citeonline{gama2026}, que estendem a abordagem para considerar indicadores de risco e
mecanismos de resposta a violações de conformidade após a implantação.

Observa-se, entretanto, que essas linhas de investigação tratam predominantemente
dimensões distintas do problema. Enquanto as métricas DORA concentram-se na avaliação
do desempenho do processo de entrega de software, os trabalhos de
\citeonline{gama2024} e \citeonline{gama2026} concentram-se na conformidade contínua
de vulnerabilidades e na resposta a alterações da postura de segurança. No conjunto de
trabalhos analisados nesta pesquisa, não foi identificado um estudo que avalie conjuntamente
uma arquitetura CI/CD baseada em Kubernetes considerando a elasticidade da infraestrutura,
o desempenho da entrega por meio de métricas DORA e mecanismos de conformidade contínua
pós-implantação fundamentados em \textit{Zero Trust}. Essa interseção constitui a lacuna
de pesquisa investigada neste trabalho.

\section{Problema de Pesquisa}
\label{sec:problema-pesquisa}

A literatura analisada evidencia duas preocupações relevantes relacionadas à
modernização do processo de entrega de software, porém investigadas predominantemente
sob perspectivas distintas. De um lado, as métricas DORA permitem avaliar quantitativamente
o desempenho da entrega de software, considerando características associadas à velocidade
e à estabilidade das mudanças \cite{dora2024}. De outro, \citeonline{gama2024} demonstra
que verificações de vulnerabilidades restritas à pipeline CI/CD não são suficientes para
lidar com alterações na postura de segurança ocorridas após a implantação.

Essa limitação decorre do caráter dinâmico das dependências de software: uma aplicação
considerada segura no momento do \textit{deploy} pode deixar de atender aos critérios de
conformidade quando novas vulnerabilidades são posteriormente identificadas. Nesse cenário,
\citeonline{gama2024} propõe associar conformidade contínua a mecanismos de identidade
baseados em \textit{Zero Trust}, de forma que \textit{workloads} que deixem de cumprir
a postura mínima de segurança possam ter sua comunicação restringida.

Considerando ambientes de entrega baseados em Kubernetes, torna-se relevante investigar
como essa dimensão de segurança pode ser analisada em conjunto com o desempenho operacional
do processo de entrega e com a elasticidade proporcionada pelo provisionamento dinâmico
de agentes de execução.

Diante desse contexto, estabelece-se a seguinte questão de pesquisa:

\begin{citacao}
Como uma arquitetura de CI/CD baseada em Kubernetes e orientada por princípios \textit{Zero Trust}, pode contribuir para o desempenho operacional e a segurança de pipelines de entrega de software?
\end{citacao}

\section{Justificativa}
\label{sec:justificativa}

A relevância deste trabalho está associada à consolidação das tecnologias \textit{Cloud
Native} como base para o desenvolvimento e a operação de sistemas modernos. Dados da
\textit{Cloud Native Computing Foundation} (CNCF) indicam que, em 2024, 89\% das
organizações pesquisadas já adotavam tecnologias \textit{Cloud Native}, enquanto 93\%
utilizavam, testavam ou avaliavam o Kubernetes. O uso da plataforma em ambientes de
produção alcançou 80\% das organizações pesquisadas, frente a 66\% no levantamento
anterior \cite{cncf2024}.

A automação dos processos de integração e entrega também acompanha essa expansão.
Segundo o mesmo levantamento, o uso de CI/CD em produção cresceu 31\% entre 2023 e
2024, sendo empregado para a maior parte ou para todas as aplicações por 60\% das
organizações pesquisadas. Além disso, 29\% dos respondentes informaram realizar
implantações de código múltiplas vezes ao dia, percentual superior aos 23\% observados
no ano anterior \cite{cncf2024}. Esses indicadores evidenciam que automação,
orquestração e ciclos frequentes de entrega deixaram de constituir práticas restritas a
ambientes experimentais e passaram a integrar infraestruturas de produção.

Entretanto, a ampla adoção dessas tecnologias não elimina os desafios relacionados ao
desempenho e à segurança do processo de entrega. O relatório DORA de 2024 identificou
quatro grupos distintos de desempenho entre as organizações pesquisadas. Apenas 19\%
dos respondentes foram classificados no nível \textit{Elite}, caracterizado por
\textit{lead time} inferior a um dia, múltiplas implantações diárias, taxa de falha de
mudanças de aproximadamente 5\% e recuperação de implantações malsucedidas em menos de
uma hora. Em contraste, 25\% foram classificados no nível de baixo desempenho, no qual
o tempo entre uma alteração e sua disponibilização pode variar entre um e seis meses e
a recuperação de uma implantação malsucedida pode demandar entre uma semana e um mês
\cite{dora2024report}.

Sob a perspectiva da segurança, o levantamento da CNCF também demonstra que a
automação das verificações ainda não acompanha integralmente a adoção da infraestrutura:
57\% das organizações pesquisadas declararam utilizar ferramentas automatizadas para
detecção de vulnerabilidades \cite{cncf2024}. Além disso, \citeonline{gama2024}
demonstra que verificações realizadas exclusivamente durante a pipeline CI/CD não
contemplam adequadamente mudanças na postura de segurança que podem ocorrer após a
implantação, como a descoberta posterior de vulnerabilidades em componentes já
executados em produção.

Nesse contexto, justifica-se investigar uma arquitetura que trate conjuntamente
eficiência operacional, desempenho da entrega e segurança contínua. A utilização de
Jenkins sobre Kubernetes com agentes efêmeros permite explorar elasticidade e isolamento
da infraestrutura de CI/CD; as métricas DORA fornecem mecanismos quantitativos para
avaliação do desempenho do processo de entrega; e os princípios de \textit{Zero Trust},
associados ao SPIRE, possibilitam incorporar mecanismos de validação de identidade e
conformidade contínua após a implantação. Assim, a pesquisa busca aproximar dimensões
que, conforme discutido nos trabalhos relacionados, são predominantemente avaliadas
separadamente na literatura analisada.

\section{Objetivos}
\label{sec:objetivos}

\subsection{Objetivo Geral}
\label{sec:objetivo-geral}

Desenvolver uma arquitetura para automação de pipelines CI/CD baseada em Kubernetes, considerando desempenho e mecanismos de segurança sob princípios de \textit{Zero Trust}.

\subsection{Objetivos Específicos}
\label{sec:objetivos-especificos}

\begin{alineas}
    \item Definir os componentes e mecanismos de uma arquitetura de CI/CD baseada em Kubernetes com provisionamento dinâmico de agentes de compilação;
    \item Implementar uma pipeline DevSecOps contendo mecanismos automatizados de análise estática, varredura de vulnerabilidades e implantação contínua;
    \item Avaliar estratégias de redução de custos utilizando segregação de workloads entre instâncias On-Demand e Spot;
\end{alineas}

\section{Organização do Trabalho}
\label{sec:organizacao-trabalho}

Este trabalho está estruturado em seis capítulos. O Capítulo 2 apresenta os fundamentos teóricos necessários para compreensão da pesquisa. O Capítulo 3 apresenta os trabalhos relacionados. O Capítulo 4 descreve a metodologia adotada e a arquitetura proposta. O Capítulo 5 apresenta os resultados obtidos e a discussão. Por fim, o Capítulo 6 apresenta as conclusões, limitações e sugestões de trabalhos futuros.
